\section{An exact rational expansion of the fixed height function}
\label{zs-section}

Retain the character, twice-normalized local heights and minimal models
of Sections~\ref{sec:height-normalization} and~\ref{sec:local-height-values}:
\begin{align}
 C&:W^2=z^6-27z^4+99z^2-9,\\
 E_1&:y^2=x^3-9x-9,&P_1&=(-2,1),\\
 E_2&:y^2=x^3-189x+999,&P_2&=(6,9),\\
 R_1&=\left(\frac{z^2-9}{4},\frac W8\right),&
 R_2&=\left(\frac{33-9/z^2}{4},-\frac{9W}{8z^3}\right).
 \label{zs-models}
\end{align}
The maps are interpreted projectively at their exceptional fibers.
Write \(\omega=dx/(2y)\), \(t=-x/y\), and \(L=\log_E\).
The unit-root complementary line is spanned by \((x+c)\omega\),
with \(c\in\mathbb Z_5\). The line spanned by \(x\omega\) is also an
allowed isotropic complement. Subscripts \(c\) and \(0\) denote their
local heights \(\lambda\), global diagonal heights \(H\), and sigma
functions. The local principal term is \(-2\log_5t\).

\begin{proposition}[Cancellation of the complementary line]
\label{zs-splitting}
For every nonzero local point, and for the infinite-order rational
points \(P_i\) above,
\begin{align}
 \lambda_c(Q)&=\lambda_0(Q)+cL(Q)^2,\\
 H_c(P_i)&=H_0(P_i)+cL(P_i)^2,&
 \alpha_c&=\alpha_0+c,\\
 \lambda_c(Q)-\alpha_cL(Q)^2
   &=\lambda_0(Q)-\alpha_0L(Q)^2,
 \label{zs-cancellation}
\end{align}
where \(\alpha_j=H_j(P_i)/L(P_i)^2\). In particular the fixed
rank-one expression does not require computing either unit-root slope.
\end{proposition}
\begin{proof}
The defining sigma differential equation gives, on the formal group,
\[
 \sigma_c(t)=\sigma_0(t)\exp(-cL(t)^2/2).
\]
Taking \(-2\log_5\) proves the first formula there. The normalized
multiplication law is
\begin{equation}
 \lambda_j([n]Q)=n^2\lambda_j(Q)-2\log_5\psi_n(Q),
 \qquad j\in\{0,c\},
 \label{zs-multiplication}
\end{equation}
for positive integers \(n\), away initially from the displayed poles.
Here \(\psi_n\) is the usual division polynomial. With our parameter,
its leading term is
\[
 \psi_n(t)=(-1)^{n-1}n\,t^{1-n^2}(1+O(t));
\]
in particular \(\psi_2=2y\sim-2t^{-3}\). The sign has logarithm zero.

For completeness, the addition identity of
\cite[Lemma~7.4]{BalakrishnanDograQCI} yields
\eqref{zs-multiplication} with the standard division-polynomial
recurrences. Its divisor differences must first be identified with
group differences: translation preserves \(\omega\), acts trivially
on \(H^1\), preserves the selected complementary line, and transports
the invariant tangent vectors. Functoriality of the local height then
identifies the height of \((Q_1)-(Q_2)\) with
\(\lambda_j(Q_1-Q_2)\). The recurrence step uses
\[
 x([n]Q)-x(Q)=-\frac{\psi_{n+1}(Q)\psi_{n-1}(Q)}{\psi_n(Q)^2}.
\]
Substitution in the addition identity gives the \(n+1\) law from
the \(n,n-1\) laws. The tangent limit supplies the doubling case;
the leading term at \(O\) fixes the constant. These conventions agree
with \cite[\S5.5 and Theorem~5.21]{SutherlandDivisionPolynomials2023}.
Continuity deals with removable exceptional charts.

Both reductions have nine points over \(\mathbb F_5\). Hence \([9]Q\)
lies in the formal group for every \(Q\in E_i(\mathbb Q_5)\).
Subtract the two versions of \eqref{zs-multiplication} with \(n=9\)
and divide by \(81\). This proves the first formula off the nonzero
nine-torsion points, and continuity extends it across them. The
heights at the other places do not change. Summing on a rational
point gives the global formula. The already established logarithms
of \(P_i\) are nonzero, so division gives the remaining identities.
No assertion that \(P_i\) generates the global group is used.
\end{proof}

\begin{proposition}[Rational recursions and the entire omitted tail]
\label{zs-tail}
For either integral pair \((a,b)=(-9,-9),(-189,999)\), let
\begin{align}
 w&=t^3+atw^2+bw^3,&w&\in t^3\mathbb Z[[t]],\\
 x&=t/w,&A(t)&=\frac{tw'(t)/w(t)-1}{2},\\
 L(t)&=\int A(t)\,dt,&
 R_0(t)&=-\int\left[\frac1t+A(t)\int x(t)A(t)\,dt\right]dt.
 \label{zs-recurrences}
\end{align}
All integrals have constant zero, including the Laurent primitive.
These are the actual invariant logarithm and
\(R_0=\log(\sigma_0/t)\). They satisfy
\[
 L\in t\mathbb Q[[t^2]],\quad L=t+O(t^5),\qquad
 R_0\in t^4\mathbb Q[[t^2]].
\]
Writing \(L=\sum l_jt^j\) and \(R_0=\sum r_jt^j\), for \(j\ge1\),
\begin{equation}
 v_5(l_j)\ge-\lfloor\log_5j\rfloor,\qquad
 v_5(r_j)\ge-2\lfloor\log_5j\rfloor.
 \label{zs-coefficient-bounds}
\end{equation}
If \(L_N,R_N\) retain degrees at most the integer \(N\ge1\), then,
uniformly for \(t\in5\mathbb Z_5\), both omitted tails have valuation
at least \(K=\lceil(N+1)/2\rceil\).
\end{proposition}
\begin{proof}
Coefficient induction gives the unique integral solution for \(w\):
each substitution strictly raises the order of an unknown correction.
It is the formal coordinate \(-1/y\). Direct differentiation gives
\(\omega=A(t)dt\), \(A=1+O(t^4)\), and
\(xA=t^{-2}+O(t^2)\). Thus the apparent \(1/t\) terms cancel.
Integrate the defining equation
\[
 -d\left(\frac{d\log\sigma_0}{\omega}\right)=x\omega.
\]
Oddness fixes the inner integration constant; removing \(\log t\)
gives precisely \eqref{zs-recurrences}, including the stated parity.

The invariant differential is integral at five, so its termwise
primitive gives the bound for \(l_j\). By
\cite[Theorem~1.3]{MazurSteinTateHeight}, the canonical odd sigma has
\(\sigma_c/t\in1+t^2\mathbb Z_5[[t]]\), and \(c\in\mathbb Z_5\).
Proposition~\ref{zs-splitting} gives
\[
 R_0=\log(\sigma_c/t)+(c/2)L^2.
\]
In degree \(j\), the logarithm denominators are integers at most
\(j/2\); this part has valuation at least
\(-\lfloor\log_5j\rfloor\). Every product contributing to \(L^2\)
has valuation at least \(-2\lfloor\log_5j\rfloor\).
This proves \eqref{zs-coefficient-bounds}, using only the existence
and integrality of the canonical sigma, not its computed coefficients.

For \(k=\lfloor\log_5j\rfloor\ge1\), induction gives
\(j\ge5^k\ge4k\). Thus
\(j-2\lfloor\log_5j\rfloor\ge\lceil j/2\rceil\); the case \(k=0\)
is immediate. Every omitted evaluated term has valuation at least
\(K\), tending to infinity with \(j\). Completeness and the
ultrametric inequality bound the full tails. The bound for \(L\)
is stronger still.
\end{proof}

\begin{proposition}[A rational-function and rational-series formula]
\label{zs-expression}
Set, as rational functions after cancellation,
\begin{align}
 \delta_i(Q)&=\frac{t([9]Q)}{\psi_9(Q)},&
 T_i&=t([9]R_i),&
 \Xi&=z^{81}\frac{\delta_1(R_1)}{\delta_2(R_2)}.
 \label{zs-rational-functions}
\end{align}
Each \(T_i\) is analytic on \(C(\mathbb Q_5)\) with values in
\(5\mathbb Z_5\), and \(\Xi\) is finite and nonzero there.
With \(\alpha_i^0=H_0(P_i)/L(P_i)^2\), the extended fixed height
function is
\begin{align}
 F_0={}&-\frac2{81}\bigl(\log_5\Xi+R_{01}(T_1)-R_{02}(T_2)\bigr)
 \nonumber\\
 &-\alpha_1^0\left(\frac{L_1(T_1)}9\right)^2
  +\alpha_2^0\left(\frac{L_2(T_2)}9\right)^2.
 \label{zs-height-formula}
\end{align}
It equals the canonical-splitting expression of
Section~\ref{sec:local-height-values}. We shall test
\(f=F_0-\Omega\), where \(\Omega=-\frac43\log_5 3\).
\end{proposition}
\begin{proof}
Near \(O\), the leading terms give
\(\delta_i=t^{81}(1+O(t))\). At a nonzero nine-torsion point,
\(t([9]Q)\) and \(\psi_9(Q)\) both vanish simply, since \([9]\)
is etale in characteristic zero; cancellation gives a nonzero value.
At any other nonzero local point, \([9]Q\) is a nonzero formal point,
where \(t\) is finite and nonzero. Thus \(\delta_i\) has no other
local zero or pole.

Away from zero and infinity, neither \(R_i\) is \(O\). At \(z=0\),
with \(W_0^2=-9\),
\(t(R_2)/z\longrightarrow-2/W_0\), so the order \(81\) in
\(\delta_2\) cancels \(z^{81}\). At infinity put \(q=1/z\) and
\(W/z^3\longrightarrow\epsilon\in\{1,-1\}\). Then
\(t(R_1)/q\longrightarrow-2/\epsilon\), cancelling the other pole.
The second quotient is finite in each respective exceptional case.
This proves all asserted local extensions, without assuming that the
zero-fiber points are rational over \(\mathbb Q\).

Off these fibers, \eqref{zs-multiplication} gives
\[
 \lambda_0(Q)=-\frac2{81}\bigl(\log_5\delta_i(Q)+R_{0i}(t([9]Q))\bigr).
\]
Subtract the two heights and \(2\log_5z\), and use
\(L(R_i)=L_i(T_i)/9\), to obtain \eqref{zs-height-formula}.
The preceding cancellations extend it across the exceptional fibers;
\eqref{zs-cancellation} identifies the two splittings.
\end{proof}

\paragraph{Exact constants and precision.}
Write the already established ninefold rational points as
\(9P_i=(A_i/d_i^2,B_i/d_i^3)\), and set
\(t_i=t(9P_i)\), \(q_i=-A_i/B_i=t_i/d_i\). The checked
identity-component hypotheses from Section~\ref{sec:height-normalization}
give
\begin{equation}
 H_0(P_i)=-\frac2{81}\bigl(\log_5q_i+R_{0i}(t_i)\bigr),
 \qquad L(P_i)=L_i(t_i)/9.
 \label{zs-constant-formula}
\end{equation}
This follows either by changing the splitting in the canonical global
formula or by its sigma definition. Each \(q_i\) is a rational
five-adic unit, whose logarithm is
\(\frac14\log(q_i^4)\). The exact leading-digit calculation gives
\(v_5H_0(P_i)=v_5L(P_i)=1\), hence \(v_5\alpha_i^0=-1\).
Changing the slope adds \(c_iL(P_i)^2\in25\mathbb Z_5\), so it does
not alter the height's first nonzero digit.

For the integer \(K=\lceil(N+1)/2\rceil\), truncating the two
\(R_0,L\) series at degree \(N\) in \eqref{zs-height-formula}
changes their combined contribution by an error in \(5^K\mathbb Z_5\),
if the constants are exact. Indeed each logarithm and its truncation
has valuation at least one; the squaring error has valuation at least
\(K+1\), and multiplication by \(\alpha_i^0\) loses one digit.
An absolute error at least \(K-2\) in \(\alpha_i^0\) also suffices,
since its multiplier is a square in \(25\mathbb Z_5\).
Evaluation errors in the rational functions and in \(\log_5\Xi\)
must be bounded separately.

The rational recurrence replay through degree forty, with
\eqref{zs-constant-formula}, yields the following exact residues:
\[
\begin{array}{c|rrr}
 &H_0(P_i)\bmod5^8&L(P_i)/5\bmod5^8&5\alpha_i^0\bmod5^8\\\hline
 i=1&245730&36459&352766\\
 i=2&96385&378472&277078
\end{array}
\]
The full-tail bounds are twenty-one absolute digits for the series
and nineteen for \(\alpha_i^0\); thus the stated residues are within
the proved precision. The smaller residues needed below are
\begin{equation}
\begin{array}{c|rrr}
 &H_0(P_i)\bmod125&L(P_i)/5\bmod25&5\alpha_i^0\bmod25\\\hline
 i=1&105&9&16\\
 i=2&10&22&3
\end{array}
 \label{zs-small-constants}
\end{equation}

\paragraph{All local charts and verification scope.}
The sextic is \(1\) at \(z=0\) modulo five and \(4\) at
\(z=\pm1,\pm2\). Each of the five residues has two nonzero simple
ordinates, and infinity has two points. Thus all local points lie
in exactly twelve smooth residue disks. Their lifts may be based at
\[
 (0,\pm\sqrt{-9}),\quad(\pm1,\pm8),\quad
 (\pm2,\pm\sqrt{19}),\quad\infty_\pm.
\]
The finite charts use \(z-a\in5\mathbb Z_5\); infinity uses
\(1/z\in5\mathbb Z_5\). The derivative \(2W\) being a unit gives
each unique analytic square-root branch. This is an inventory of
local disks, not of rational points.

The exact rational recurrences and residues are replayed by
\texttt{next\_replay\_zero\_slope.py}; the canonical result and exact
source hashes are listed in the accompanying geometry inventory.
The entire-tail estimates, splitting comparison and projective
extensions above are ordinary proofs using the cited primary inputs.
They are not assertions of Lean formalization or consequences of
finite coefficient inspection. The following disk certificates and
the separate rational sieve use these estimates to complete the
fixed-curve calculation; none supplies a uniform height bound for
varying exponents or residuals.
